% ============================================================
% Section 2 — Related Work
% ============================================================
\section{Related Work}
\label{sec:related_work}

\subsection{Continual Learning for Language Models}

Continual learning (CL) aims to update model parameters on a stream of tasks without catastrophic forgetting~\citep{kirkpatrick2017overcoming,zenke2017continual}.
Regularisation-based methods penalise changes to important weights~\citep{kirkpatrick2017overcoming,li2017learning}, while replay-based methods store or generate past examples for interleaved training~\citep{de2021continual,scialom2022fine}.
Recent work has extended CL to large language models, studying instruction tuning across sequential tasks~\citep{wu2024continual} and domain-incremental adaptation~\citep{ke2023continual}.
However, existing CL methods assume a \emph{given} data stream and do not address the problem of \emph{which} experiences should be retained.
HAT complements CL by introducing a principled experience selection mechanism upstream of parameter updates.

\subsection{Memory-Augmented Neural Networks}

Memory-augmented architectures equip neural networks with external memory stores that can be read from and written to~\citep{graves2014neural,sukhbaatar2015end}.
More recent work combines retrieval mechanisms with language models~\citep{borgeaud2022improving,zhong2022training}, enabling access to a large knowledge base at inference time.
While these approaches augment \emph{inference}, they do not leverage memory for \emph{learning}: retrieved information is used contextually but does not update the model's parameters.
In contrast, HAT uses its memory store (the Neocortex) as a curated training set that feeds back into parameter optimisation.

\subsection{Retrieval-Augmented Generation}

Retrieval-augmented generation (RAG) methods prepend retrieved passages to the model's input context~\citep{lewis2020retrieval,guu2020realm,izacard2022atlas}.
RAG is effective for knowledge-intensive tasks but inherits the limitations of fixed-context approaches: the model cannot \emph{learn} from retrieved content, context windows are bounded, and retrieval quality degrades when the knowledge base is noisy.
HAT can be seen as going beyond RAG by closing the loop from retrieval to training: high-value interactions—including oracle-provided corrections—are not only used at inference time but are consolidated into the model's parameters.

\subsection{Knowledge Distillation and Self-Improvement}

Knowledge distillation transfers knowledge from a teacher model to a student~\citep{hinton2015distilling,sanh2019distilbert}.
Recent work explores \emph{self-improving} language models that generate their own training data~\citep{huang2023large,yuan2024self,chen2024self} or use reinforcement learning from human/AI feedback~\citep{ouyang2022training,bai2022constitutional}.
These methods typically apply distillation or self-play uniformly, without selecting \emph{which} experiences are most informative.
HAT's Oracle consultation mechanism is a form of selective, on-demand distillation: the teacher is queried only when the student model is uncertain or has been corrected by the user, ensuring efficient use of the oracle's capacity.

\subsection{Complementary Learning Systems Theory}

Complementary learning systems (CLS) theory posits that the brain uses two interacting memory systems: the hippocampus for rapid encoding of specific episodes and the neocortex for slow extraction of statistical regularities~\citep{mcclelland1995there,kumaran2016learning}.
Sleep replay, particularly during slow-wave sleep, is hypothesised to mediate the transfer from hippocampus to neocortex~\citep{rasch2013sleep,walker2017sleep}.
Several works in machine learning have drawn on CLS for continual learning~\citep{shin2017continual,van2020brain,arani2022learning}, but they focus on image classification and do not address the unique challenges of language model self-improvement from interactive experience.
HAT directly operationalises the CLS framework for LLMs, with an explicit hippocampal selection stage, a neocortical memory store, and an offline SWS training phase.
